\section{Games and Boolean functions}
\label{sec:preliminaries}

For a finite set \(S\), we identify \(\{0,1\}^{S}\) with the power set
\(2^{S}\).  For \(p\in(0,1)\), let \(\mu_p^{S}\) be the probability
measure on \(2^{S}\) in which each element belongs to the sampled set
independently with probability \(p\).  We use \([n]=\{1,\ldots,n\}\).
All graphs and probability spaces in the reduction are finite.

\subsection{The game used in the reduction}

A two-to-one game consists of a bipartite constraint graph with disjoint
vertex sets \(U,V\), nonempty label sets \(X,Y\) with
\(\abs{X}=2\abs{Y}\), and, for each edge \(e=(u,v)\), a map
\(\pi_e:X\to Y\) for which every fiber has cardinality two.  Parallel
constraints are allowed.  The edges carry a rational probability
distribution \(\nu\).  Its endpoint marginals are denoted by
\(\nu_U\) and \(\nu_V\).  A labeling consists of maps
\(\ell_U:U\to X\) and \(\ell_V:V\to Y\); its value is
\[
  \Pr_{e=(u,v)\sim\nu}
  [\,\pi_e(\ell_U(u))=\ell_V(v)\,].
\]
The value of a game is the maximum of this expression over all
labelings.

\begin{theorem}[Two-to-one games with imperfect completeness]
\label{thm:games}
For every constant \(\theta\in(0,1/2)\), there are constant label
sets \(X,Y\) such that it is \(\mathrm{NP}\)-hard, by a deterministic
polynomial-time reduction, to distinguish two-to-one games of the
following kinds:
\begin{enumerate}[label=\textup{(\roman*)}]
  \item there is a labeling of value at least \(1-\theta\);
  \item every labeling has value at most \(\theta\).
\end{enumerate}
Every constraint is a total map \(X\to Y\) with exactly two preimages
of each label in \(Y\).  The label sets may depend on \(\theta\).
\end{theorem}

We invoke the completed theorem in this precise form.
Dinur, Khot, Kindler, Minzer, and Safra prove the stated gap
conditionally on their Grassmann agreement hypothesis
\cite[Theorem~1.8]{DKKMS2025}.  Their Hypothesis~3.6 follows from the
Grassmann expansion theorem of Khot, Minzer, and Safra
\cite[Theorem~1.12, ECCC revision~2]{KMS2018} through the implications
of Barak, Kothari, and Steurer
\cite[Section~1.1 and Theorems~9--11]{BKS2019}.
The two-to-one definition and the common alphabets are as in
\cite[Definitions~1.1--1.2 and Section~4.2]{DKKMS2025}.
In the explicit constraint construction, the smoothing parameter may
be chosen rational within its permitted range
\cite[Sections~4.2 and~5.2]{DKKMS2025}. The edge sampling process then
specifies rational weights: all constant-length tuples and constant-dimensional
subspaces can be enumerated, followed by the prescribed identifications.
With the parameters fixed, this is a deterministic finite construction
of polynomial size.  Thus the use of a rational edge law does not
introduce randomness into the reduction.  We require neither perfect
completeness nor any degree regularity of the game.

\subsection{Junta approximation at fixed dyadic biases}

For a Boolean function \(f:2^S\to\{0,1\}\) and \(j\in S\), define
\begin{equation}
\label{eq:biased-influence}
  I_{p,j}(f)
  =\Pr_{R\sim\mu_p^{S\setminus\{j\}}}
       [\,f(R)\ne f(R\cup\{j\})\,],
  \qquad
  I_p(f)=\sum_{j\in S} I_{p,j}(f).
\end{equation}
This is also the probability that flipping coordinate \(j\) changes
the value of \(f\), when the entire input has law \(\mu_p^S\):
the change event depends only on the other coordinates.
A function is a \emph{junta on \(J\subseteq S\)} if its value depends
only on \(R\cap J\).  Constant functions are juntas on the empty set.

\begin{theorem}[Friedgut's junta theorem]
\label{thm:friedgut}
For every \(b\ge0\) and \(\alpha>0\), there is a finite integer
\(J_{\mathrm{unif}}(b,\alpha)\) with the following property.
For every finite \(S\) and every Boolean function
\(f:2^S\to\{0,1\}\) with \(I_{1/2}(f)\le b\), there is a Boolean
function \(\widetilde f\), depending on at most
\(J_{\mathrm{unif}}(b,\alpha)\) coordinates, for which
\[
  \Pr_{R\sim\mu_{1/2}^{S}}[\,f(R)\ne\widetilde f(R)\,]\le\alpha.
\]
The bound is independent of \(\abs S\).
\end{theorem}

This is the uniform version of Friedgut's theorem
\cite[Theorem~1.1 in the February 3, 1998 author manuscript]{Friedgut1998}.
Friedgut also proves a biased-product version
\cite[Theorem~4.1 in the same manuscript]{Friedgut1998}.
We give an elementary reduction of the dyadic case to the uniform case,
so that the dependence on the bias and independence of the game
alphabet size are explicit.

\begin{corollary}[Dyadic junta approximation]
\label{cor:dyadic-junta}
Let \(p=2^{-k}\), where \(k\ge1\) is an integer.  For every \(b\ge0\)
and \(\alpha>0\), there is a finite integer \(J(p,b,\alpha)\),
independent of \(\abs S\), such that every Boolean function
\(f:2^S\to\{0,1\}\) with \(I_p(f)\le b\) has a Boolean junta
approximant on at most \(J(p,b,\alpha)\) coordinates with error at
most \(\alpha\) under \(\mu_p^S\).
\end{corollary}

\begin{proof}
The assertion is immediate for \(S=\varnothing\), so assume otherwise.
Let \((X_{j,a})_{j\in S,\,a\in[k]}\) be independent uniform bits and
put
\[
  Z_j=\bigwedge_{a=1}^k X_{j,a},
  \qquad
  F(X)=f(\{j\in S:Z_j=1\}).
\]
The \(Z_j\) are independent Bernoulli \(p\) bits.  Flipping \(X_{j,a}\)
changes \(F\) precisely when all the other \(k-1\) bits of its block
are one and coordinate \(j\) is pivotal for the values of the other
blocks.  These two events are independent.  Therefore
\begin{equation}
\label{eq:and-block-influence}
  I_{1/2}(F)
   =k2^{-(k-1)}I_p(f)
   =2kpI_p(f)
   \le kb.
\end{equation}
By Theorem~\ref{thm:friedgut}, there is a Boolean approximant \(H\)
to \(F\) with error at most \(\alpha\), using at most
\(J_{\mathrm{unif}}(kb,\alpha)\) of the individual bits \(X_{j,a}\).
Let \(J\subseteq S\) be the set of blocks touched by these bits.
Then \(\abs J\le J_{\mathrm{unif}}(kb,\alpha)\).

Reveal all bits in every block indexed by \(J\).  Since \(H\) is
measurable with respect to these revealed blocks, the Boolean
predictor with smallest conditional error from the revealed blocks
has total error at most that of \(H\).  More explicitly, this
predictor outputs one when the conditional probability of \(F=1\)
is at least \(1/2\), with any fixed convention at equality.
The unrevealed blocks remain independent of the revealed blocks,
and the dependence of \(F\) on a revealed block is only through its
AND value.  Consequently
\[
 \Pr[\,F=1\mid (X_{j,a})_{j\in J,\,a\in[k]}\,]
 =
 \Pr[\,f(\{j:Z_j=1\})=1\mid (Z_j)_{j\in J}\,].
\]
The optimal predictor is therefore a Boolean function only of
\((Z_j)_{j\in J}\).  Regard it as a function
\(\widetilde f:2^S\to\{0,1\}\) on the original coordinates.
The induced law of the original input is \(\mu_p^S\), so
\[
  \Pr_{R\sim\mu_p^S}[\,f(R)\ne\widetilde f(R)\,]\le\alpha.
\]
Thus \(J(p,b,\alpha)=J_{\mathrm{unif}}(kb,\alpha)\) is a valid
choice, increased to one if desired.
\end{proof}

\begin{remark}
The junta approximants need not be monotone.  In the reduction, the
original conditioned functions will be monotone, but the argument
uses only their approximation errors and their coordinate supports.
For any fixed finite collection of dyadic biases, the maximum of the
bounds in Corollary~\ref{cor:dyadic-junta} remains independent of the
alphabet size.
\end{remark}
